\chapter{ATTITUDE AND HEADING}
\label{ch:attitude}

The robot drives on flat floors, so the navigation state needs only the yaw
$\psi$ from the IMU. Roll and pitch matter only for tilt-compensating the
magnetometer. This chapter derives the gyroscope heading used by the odometry,
its error growth, the accelerometer tilt, and the magnetometer compass.

\section{SENSOR FRAME}

\texttt{librobotcontrol} returns the gyroscope, accelerometer and
magnetometer in one common frame: the AK8963 magnetometer inside the MPU-9250
has its $x$ and $y$ axes swapped and $z$ reversed relative to the gyroscope,
and the library undoes that (\texttt{mpu.c}, $m_x\leftarrow m_y'$,
$m_y\leftarrow m_x'$, $m_z\leftarrow -m_z'$). The static transform
\frm{base\_link}$\to$\frm{imu\_link} has zero rotation, which assumes the
cape's IMU axes coincide with the robot's $x$-forward, $y$-left, $z$-up axes.
The yaw-rate sign has been confirmed in practice (gyro heading turns the right
way in RViz and agrees with the compass over two turns); the $x$ and $y$ axes
have not been checked separately.

\section{SENSOR ERROR MODEL}

Following the thesis (thesis \S7.5, after Gebre-Egziabher
\cite{gebreegziabher2004design} and Flenniken \cite{flenniken2005modeling}),
each inertial channel is modelled as
\begin{equation}\label{eq:sensormodel}
m = m_r + c_r + b_r + w_r,
\end{equation}
where $m_r$ is the true value, $c_r$ a constant bias that changes at every
power-up, $b_r$ a walking bias and $w_r$ white noise of variance $\sigma_r^2$.
The walking bias is a first-order Gauss--Markov process,
\begin{equation}\label{eq:walkingbias}
\dot b_r = -\frac{1}{\tau_r}b_r + w_{b_r},
\end{equation}
with $w_{b_r}$ white of intensity $2\sigma_{b_r}^2/\tau_r$, so that
$\sigma_{b_r}^2$ is its steady-state variance. The thesis characterised an
MPU-class gyroscope this way, with Allan variance, and found
$\tau_r\approx\SI{1600}{s}$ (thesis Table~7.2); the same procedure can be
reused on this IMU by logging \topic{/imu/data\_raw} at rest for an hour.

\section{GYROSCOPE HEADING}
\label{sec:gyro}

\subsection{Bias removal at start-up}

At start the base node averages the $z$ rate for $T_b=\SI{2}{s}$ with the
robot stationary,
\begin{equation}\label{eq:gyrobias}
\hat b_g = \frac{1}{n}\sum_{i=1}^{n} g_{z,i},
\end{equation}
which estimates $c_r+b_r(0)$ in \eqref{eq:sensormodel}. This is the same
start-up bias removal used on the TWIP (thesis \S4.3). The corrected rate
$g_z-\hat b_g$ is integrated in \eqref{eq:dpsi}.

\subsection{Heading error growth}

The heading error after time $t$ has three parts. Integrating white noise
gives an angle random walk; an error $\delta b=\hat b_g-(c_r+b_r(0))$ in the
start-up estimate gives a linear ramp; and the walking bias adds a slower
term,
\begin{equation}\label{eq:headingerr}
\delta\psi(t) = \underbrace{\int_0^t w_r\,\dd s}_{\sigma=N_g\sqrt t}
\;+\;\underbrace{\delta b\,t}_{\text{start-up error}}
\;+\;\int_0^t\bigl(b_r(s)-b_r(0)\bigr)\dd s .
\end{equation}
The MPU-9250 datasheet gives a rate-noise density of
$N_g=\SI{0.01}{\degree/s/\sqrt{Hz}}$. The random-walk term is therefore
$\SI{0.01}{\degree}\sqrt{t/\mathrm{s}}$, under \SI{0.1}{\degree} after one
minute. Averaging over $T_b=\SI{2}{s}$ leaves $\delta b$ with a standard
deviation of about $N_g/\sqrt{T_b}=\SI{0.007}{\degree/s}$, which would give
about \SI{0.4}{\degree/min}.

The measured drift at rest was \SI{1.6}{\degree/min}
(\SI{0.027}{\degree/s}), four times larger. The noise terms cannot account
for it, so most of it is bias that moves after the start-up average:
temperature change as the board and motor drivers warm up (the datasheet
allows the zero-rate offset to move by up to \SI{\pm30}{\degree/s} over the
full temperature range) and the walking bias. Over a 10-minute mapping run
this is about \SI{16}{\degree} of heading error if nothing corrects it. Two
practical consequences:
\begin{itemize}
\item Measuring the bias after the robot has been powered for a few minutes
  will reduce the error.
\item The bias should keep being estimated while the robot runs. The base
  node now does this whenever the robot is parked (\S\ref{sec:zupt}); the
  magnetometer can also bound the heading (\S\ref{sec:compfilt}).
\end{itemize}

\subsection{Zero-velocity update}
\label{sec:zupt}

When the robot is parked its true yaw rate is zero, so whatever the gyroscope
reads is bias. The base node uses this. Let $t_m$ be the last time any wheel
tick changed or a \topic{/cmd\_vel} command was active. If
$t-t_m>T_z=\SI{0.5}{s}$ the robot is declared stationary and, each cycle,
\begin{equation}\label{eq:zupt}
\hat b_{g,k} = \hat b_{g,k-1} + \alpha\bigl(g_{z,k}-\hat b_{g,k-1}\bigr),
\qquad \alpha=\min\!\left(1,\frac{\Delta t}{\tau_z}\right),
\qquad \Delta\psi_k = 0,
\end{equation}
with $\tau_z=\SI{20}{s}$. The first part is a first-order low-pass filter
(time constant $\tau_z$) that tracks the walking bias $b_r$ of
\eqref{eq:walkingbias}; the second freezes the heading. In Kalman filter terms
this is a pseudo-measurement $\omega=0$ with zero variance on the heading and a
steady-state gain $\alpha$ on the bias \cite{groves2013principles}. After the
change, heading drift while parked fell from about \SI{1.5}{\degree/min} to
zero (2 October 2026). Drift while driving is still governed by
\eqref{eq:headingerr}, but every stop now re-estimates the bias.

\begin{remark}
The stationary test uses the wheels and the command, not the gyroscope. If the
robot is picked up and turned by hand, or pushed sideways with the wheels
locked, the heading will not follow, and the bias filter will absorb part of
the real rotation. Carry the robot level and still, or restart
\texttt{bbai\_base}, when moving it by hand.
\end{remark}

\section{TILT FROM THE ACCELEROMETER}

The accelerometer measures specific force, which at rest is gravity in body
axes. With roll $\phi$ and pitch $\theta$, the thesis derives (thesis
eq.~4.2, after Kuipers \cite{kuipers1999quaternions})
\begin{equation}\label{eq:rollpitch}
\phi=\atantwo(a_y,a_z),\qquad
\theta=\atantwo\!\left(-a_x,\sqrt{a_y^2+a_z^2}\right).
\end{equation}
This assumes no acceleration other than gravity, which is close enough for
this robot's low accelerations, and could be filtered with the complementary
filter of the thesis (thesis eq.~4.4) if needed. The current software does
not use roll and pitch.

\section{MAGNETOMETER COMPASS}
\label{sec:mag}

\subsection{Measurement model}

The magnetometer measures Earth's field $\bm h$ rotated into the body frame,
distorted by nearby material, plus noise \cite{gebreegziabher2006mag}:
\begin{equation}\label{eq:magmodel}
\bm m = S\,R\T\bm h + \bm b_{\mathrm{hi}} + \bm w_m .
\end{equation}
$R$ is the body-to-world rotation. The hard-iron offset $\bm b_{\mathrm{hi}}$
is a constant field carried with the robot: permanent magnets in the motors
and encoder rings, and steady currents. The matrix $S$ combines the sensor's
scale factors, axis misalignment and soft-iron effects, where nearby steel
bends Earth's field. On this robot the measured field magnitude was about
\SI{150}{\micro T} against Earth's \SI{50}{\micro T}, so
$\bm b_{\mathrm{hi}}$ dominates.

\subsection{Planar calibration}

When the robot spins on a level floor, $R$ is a rotation about $z$ alone and
the horizontal components $(m_x,m_y)$ trace an ellipse; without soft-iron
distortion it is a circle of radius $H$ (the horizontal field) centred on
$(b_x,b_y)$. The calibration routine (\texttt{mag\_calibrate.py}) commands
\SI{5}{rad/s} until the gyroscope has seen two full turns (or \SI{40}{s} pass),
logging $(m_x,m_y)$ from \topic{/imu/mag}. A lower command stalls on the floor
(\S\ref{sec:skidturn}).

\paragraph{Circle fit (used).} The routine fits a circle by linear least
squares (the K{\aa}sa fit \cite{kasa1976circle}): writing the circle as
$m_x^2+m_y^2+Dm_x+Em_y+F=0$, each sample gives one linear equation in
$(D,E,F)$, and the normal equations
\begin{equation}\label{eq:kasa}
\begin{bmatrix}\sum m_x^2&\sum m_xm_y&\sum m_x\\
\sum m_xm_y&\sum m_y^2&\sum m_y\\
\sum m_x&\sum m_y&n\end{bmatrix}
\begin{bmatrix}D\\E\\F\end{bmatrix}
=-\begin{bmatrix}\sum m_x z\\ \sum m_y z\\ \sum z\end{bmatrix},
\qquad z=m_x^2+m_y^2,
\end{equation}
are solved by Cramer's rule. The hard-iron offset is the centre
$(o_x,o_y)=(-D/2,-E/2)$ and the horizontal field is
$H=\sqrt{o_x^2+o_y^2-F}$. The routine also unwraps the compass heading over
the spin and compares its direction with the gyroscope's to set
\texttt{mag\_sign} $\in\{+1,-1\}$. The result on 1 October 2026 was
$(o_x,o_y)=(11.71,\,66.61)$~\si{\micro T}, $s=+1$. A circle fit removes the
hard iron only; if the residual (printed as ``fit error'') grows, one of the
estimators below adds soft-iron correction.

\paragraph{Min--max (offset and per-axis scale).}
\begin{equation}\label{eq:minmax}
b_x=\tfrac12(m_x^{\max}+m_x^{\min}),\quad
s_x=\frac{\bar\rho}{\tfrac12(m_x^{\max}-m_x^{\min})},
\end{equation}
and likewise for $y$, with $\bar\rho$ the mean of the two half-ranges. The
corrected field is $m_x'=s_x(m_x-b_x)$, $m_y'=s_y(m_y-b_y)$. This handles
hard iron and unequal axis gains but not rotation of the ellipse.

\paragraph{Ellipse fit (full planar soft iron).} Fit the conic
\begin{equation}
A m_x^2 + B m_xm_y + C m_y^2 + D m_x + E m_y = 1
\end{equation}
to all samples by linear least squares. The centre
$\bm b = -\tfrac12 M^{-1}[D\ E]\T$ with
$M=\begin{bmatrix}A&B/2\\B/2&C\end{bmatrix}$ is the hard-iron offset, and the
correction $\bm m' = M^{1/2}(\bm m-\bm b)$ maps the ellipse back to a circle.

The routine prints the resulting lines for \file{base.yaml}. The check during
calibration was that over the two-turn spin the compass heading advanced
\SI{745}{\degree} while the gyroscope advanced \SI{724}{\degree}, against
\SI{720}{\degree} actual. The gyroscope's scale factor is therefore good to
about 0.6\% and the compass's residual distortion is about 3.5\%.

\subsection{Heading from the field}

On level ground, write the world (ENU) field's horizontal part in terms of the
magnetic declination $D$ (positive when magnetic north is east of true north):
$\bm h_{\mathrm{hor}} = H[\sin D,\ \cos D]\T$ in (east, north). With the robot
at ENU yaw $\psi$, the body-frame components are
\begin{equation}
m_x' = H\sin(\psi+D),\qquad m_y' = H\cos(\psi+D),
\end{equation}
so
\begin{equation}\label{eq:magyaw}
\psi = \atantwo(m_x',\,m_y') - D .
\end{equation}
This is exactly what \texttt{base\_node.py} computes, with
$m'=s\,(m-o)$ from the circle fit:
$\psi_m = s\,\atantwo(m_x-o_x,\,m_y-o_y)-D$. It publishes $\psi_m$ as the
orientation of \topic{/imu/data} with variance \SI{0.05}{rad^2}
($\sigma\approx\SI{13}{\degree}$) on yaw and $10^3$ on roll and pitch (unknown),
and the gyroscope rate with variance \SI{1e-4}{rad^2/s^2}.
Checking the two easy cases: facing east ($\psi=0$) with $D=0$, north is on
the robot's left, so $m_y'=H$, $m_x'=0$ and $\psi=0$; facing north
($\psi=\pi/2$), north is straight ahead, $m_x'=H$ and $\psi=\pi/2$.

Two conventions are easy to confuse here. A compass heading $\chi_c$ is measured
clockwise from north, while ROS yaw is counter-clockwise from east:
\begin{equation}
\chi_c = \tfrac{\pi}{2}-\psi .
\end{equation}
The robot's declination is set to \SI{2.5}{\degree} east, computed from the
NOAA World Magnetic Model \cite{wmm2025} for the robot's home location. (The
coordinates themselves are deliberately not recorded in this document or the
repository.)

\subsection{Tilt compensation}

If the robot is not level, the levelled field is
\begin{equation}
\bm m_h = R_y(\theta)R_x(\phi)\,\bm m',
\end{equation}
using the ZYX convention in which $R=R_z(\psi)R_y(\theta)R_x(\phi)$, and
\eqref{eq:magyaw} is applied to $(m_{h,x},m_{h,y})$. With $\phi$ and $\theta$
from \eqref{eq:rollpitch}. On a level floor $\bm m_h=\bm m'$. A \SI{5}{\degree}
tilt with a magnetic inclination of about \SI{60}{\degree} (typical of the
southern United States) produces up to about \SI{9}{\degree} of heading error,
so compensation matters on ramps and lawns but not indoors.

\subsection{Limits indoors}

Indoors, steel in floors, appliances and wiring adds local fields that move
the compass by tens of degrees as the robot drives. The calibration above
removes only fields carried with the robot. The compass is therefore useful
outdoors for initialising the GPS frame (\S\ref{sec:navsat}) but should not
be trusted as an absolute heading inside the house.

\section{FUSING GYRO AND COMPASS}
\label{sec:compfilt}

The gyroscope is precise over seconds and drifts over minutes; the compass is
noisy and locally biased but does not drift. The thesis used exactly this
structure for tilt (thesis eq.~4.4), and the heading version is
\begin{equation}\label{eq:yawcf}
\psi_{k+1} = \psi_k + (g_{z,k}-\hat b_g)\Delta t
  + (1-\lambda)\,\wrap\!\bigl(\psi_{m,k}-\psi_k-(g_{z,k}-\hat b_g)\Delta t\bigr),
\end{equation}
where $\psi_m$ is \eqref{eq:magyaw}. The wrap is needed because the two
headings can differ across the $\pm\pi$ boundary. The filter has a crossover
time constant $\tau=\lambda\Delta t/(1-\lambda)$: with $\Delta t=$1/30$~s$,
$\lambda=0.999$ gives $\tau\approx\SI{33}{s}$. Higgins \cite{higgins1975comparison}
shows this is the steady-state Kalman filter for a random-walk heading with
white measurement noise, which is the form the EKF of
Chapter~\ref{ch:fusion} implements with time-varying gains.
\eqref{eq:yawcf} is not currently running. The compass feeds only the GPS
EKF on the laptop; the odometry heading that gmapping uses is gyroscope only,
which is deliberate given the indoor limits above.
